\chapter{A brief excursion into measure theory}

Here we briefly mention some of the basic definitions
and results from measure theory, and point out how we
used them in the previous lectures. The relevant
material is covered in Appendix A.1 in [Dur2010]
(Appendix A.1 and A.2 in [Dur2004]). It is not
required reading, but if you read it you are perhaps
more likely to attain a good understanding of the
material that \emph{is} required...

\begin{defn}

\begin{enumerate}[label=(\roman*)]
\item A $\pi$-system is a collection ${\cal P}$ of subsets
of a set $\Omega$ that is closed under intersection
of two sets, i.e., if $A,B\in {\cal P}$ then
$A\cap B\in {\cal P}$.

\item A $\lambda$-system is a collection ${\cal L}$ of
subsets of a set $\Omega$ such that: 1.
$\Omega \in {\cal L}$; 2. If $A,B\in {\cal L}$ and
$A\subset B$ then $B\setminus A\in {\cal L}$; 3. If
$(A_n)_{n=1}^\infty$ are all in ${\cal L}$ and
$A_n \uparrow A$ then $A\in {\cal L}$.
\end{enumerate}

\end{defn}

\noindent The following is a somewhat technical result that
turns out to be quite useful:

\begin{thm}[Dynkin's $\pi-\lambda$ theorem]
If ${\cal P}$ is a $\pi$-system and ${\cal L}$ is a
$\lambda$-system that contains ${\cal P}$ then
$\sigma({\cal P}) \subset {\cal L}$.
\end{thm}

\begin{lem}[Uniqueness theorem]
If the values of two probability measures $\mu_1$ and
$\mu_2$ coincide on a collection ${\cal P}$ of sets,
and ${\cal P}$ is a $\pi$-system (closed under finite
intersection), then $\mu_1$ and $\mu_2$ coincide on
the generated $\sigma$-algebra $\sigma({\cal P})$.
\end{lem}

\noindent The uniqueness theorem implies for example that to
check if random variables $X$ and $Y$ are equal in
distribution, it is enough to check that they have the
same distribution functions.

Both of the above results are used in the proof of
the following important theorem in measure theory:

\begin{thm}[Carathéodory's extension theorem]
Let $\mu$ be an almost-probability-measure defined on
an algebra ${\cal A}$ of subsets of a set $\Omega$.
That is, it satisfies all the axioms of a probability
measure except it is defined on an algebra and not a
$\sigma$-algebra; $\sigma$-additivity is satisfied
whenever the countable union of disjoint sets in the
algebra is also an element of the algebra. Then $\mu$
has a unique extension to a probability measure on the
$\sigma$-algebra $\sigma({\cal A})$ generated by
${\cal A}$.
\end{thm}

\noindent Carathéodory's extension theorem is the main tool
used in measure theory for constructing measures: one
always starts out by defining the measure on some
relatively small family of sets and then extending to
the generated $\sigma$-algebra (after verifying
$\sigma$-additivity, which often requires using
topological arguments, e.g., involving compactness).
Applications include:

\begin{itemize}
\item Existence and uniqueness of Lebesgue measure in
$(0,1)$, $\R$ and $\R^d$.

\item Existence and uniqueness of probability measures
associated with a given distribution function in $\R$
(sometimes called ``Lebesgue-Stieltjes measures in
$\R$''). We proved existence instead by starting with
Lebesgue measure and using quantile functions.

\item Existence and uniqueness of Lebesgue-Stieltjes
measures in $\R^d$ (i.e., measures associated with a
$d$-dimensional joint distribution function). Here
there is no good concept analogous to quantile
functions, although there are other ways to construct
such measures explicitly using ordinary Lebesgue
measures.

\item Product measures -- this corresponds to probabilistic
experiments which consist of several independent
smaller experiments.
\end{itemize}

\noindent Note that Durrett's book also talks about measures
that are not probability measures, i.e., the total
measure of the space is not 1 and may even be
infinite. In this setting, the theorems above can be
formulated in greater generality.
